%% ma: L12-B08-0042
%% lop: 12
%% chuong: 2
%% bai: 8
%% dang: 12.8.30
%% loai: TN4
%% muc: VD
%% dap_an: "A"
%% nguon: "LaTeX giáo viên gửi 21/10/2026 (M05, Vectơ và hệ trục tọa độ trong không gian, 107 câu), Phần I Câu 35"
%% kien_thuc: "Bài 8 (khoảng cách hai điểm), lớp 10 (giá trị nhỏ nhất của tam thức bậc hai)"
%% trang_thai: chinh-thuc
%% ghi_chu: "Giáo viên đã duyệt 21/10/2026. Phần I Câu 50 của tài liệu cùng mô hình, chỉ đổi số và đổi cách hỏi (hỏi tọa độ M) nên đã bỏ"
\begin{ex}
Trong không gian $Oxyz$, cho hai điểm $A(1;2;3)$ và $B(-1;4;1)$. Gọi $M(x;y;0)$ là điểm di động trên mặt phẳng $(Oxy)$ sao cho biểu thức $P=MA^2+MB^2$ đạt giá trị nhỏ nhất. Tổng $S=x+y$ bằng
\choice{\True $3$}{$0$}{$1$}{$-2$}
\loigiai{$P=(x-1)^2+(y-2)^2+9+(x+1)^2+(y-4)^2+1=2x^2+2(y-3)^2+16\ge16$.

Dấu bằng khi $x=0$, $y=3$. Vậy $S=3$.}
\end{ex}
